\documentclass[
  paper=a4,
  fontsize=12pt,
  DIV=16,
  headheight=52pt,
  footheight=45pt,
  headinclude,
  parskip=full,
]{scrartcl}

\usepackage{fontspec}
\setsansfont{Source Sans Pro}
\renewcommand{\familydefault}{\sfdefault}

\usepackage{polyglossia}
\setmainlanguage{german}

\usepackage{csquotes}

\usepackage{microtype}

\usepackage{scrlayer-scrpage}
\pagestyle{scrheadings}
\setkomafont{pagehead}{\normalfont}
\setkomafont{pagefoot}{\normalfont\footnotesize}
\ohead{\includegraphics[height=1.5cm]{pep.pdf}}
\ihead{%
  \large\bfseries PeP et al.\ e.\,V.\\%
  Begabtenförderungsprogramm –\\
  Deutschlandstipendien%
}
\cfoot{}
\ifoot{%
  PeP et al.~e.\,V.\\
  Begabtenförderung\\
  \url{www.pep-dortmund.org}
}


\usepackage{titling}
\usepackage{booktabs}

\usepackage{xcolor}
\usepackage{calc}
\usepackage[colorlinks=true,urlcolor=blue!50!black]{hyperref}

\date{31. Juli 2014}

\begin{document}
\textbf{\Huge\sffamily Stipendium für Physikstudierende}\\[0.5\baselineskip]
\textbf{\Large\sffamily Studienförderung für  Physikerinnen in Dortmund}

{\large Der Verein der Physikstudierenden und ehemaligen Physikstudierenden \enquote{PEP et al. e.\,V.}
vergibt auch im Wintersemester 2015/2016 wieder Stipendien an Dortmunder Physikerinnen und Physiker.
Dieses von der Bundesregierung unterstütze Programm wird zur
Hälfte aus Vereinsmitteln und zur anderen Hälfte aus dem Bundeshaushalt finanziert.
}

Das Programm richtet sich an Studierende mit überdurchschnittlichen Studienleistungen und dient explizit der Förderung des wissenschaftlichen Nachwuchses der Dortmunder Fakultät.

Die finanzielle Förderung beträgt 300 € monatlich und ist auf ein Jahr befristet.
Die Vergabe der Stipendien erfolgt leistungsorientiert und einkommensunabhängig.
Es werden fünf Stipendien vergeben.

Eine Bewerbung ist möglich, wenn Du an der Technischen Universität Dortmund als Studierender eingeschrieben bist und Physik im Studiengang Bachelor of Science, Bachelor
Lehramt Physik (BfP) oder Medizinphysik studierst und zum Wintersemester im dritten
oder fünften Fachsemester bist.

Das Auswahlverfahren beinhaltet eine schriftliche Bewerbung und eine anschließende
Beurteilung durch ein Gremium aus Hochschullehrern und Mitgliedern von PEP et al. e.\,V.
Eine Bewerbung ist jederzeit, spätestens jedoch bis zum 31. Juli 2014 möglich.
Alle nötigen Informationen und Dokumente erhältst Du auf der Internetseite des Vereins

\begin{center}
  \Huge\bfseries\url{http://stipendium.pep-dortmund.org}
\end{center}

Für persönliche Fragen stehen Dir Julian Wishahi und Christophe Cauet (P2-O1-424) zur
Verfügung.
\end{document}
